\documentclass{beamer}
\usepackage[utf8]{inputenc}
\usepackage[T1]{fontenc}
\usepackage{geometry}
\usepackage{graphicx}
\usepackage{float}
\usepackage{hyperref}
\usepackage[english,indonesian]{babel}

\title{Optimasi Pemrosesan Template LaTeX Menggunakan Pendekatan Berbasis AI}
\author{Academic Support Agent}
\date{\today}

\begin{document}
\maketitle

\begin{abstract}

\end{abstract}

\chapter{PENDAHULUAN}


\chapter{TINJAUAN PUSTAKA}


% Original content (commented out):
% %颜色没有大调整，跟校徽色不符
%只在参考的两个模板上改动了颜色，以及解决了参考文献的问题。
%[1]中所附的一些排版内容在本模板中删去，有需要可以从[1][2]中学习。

\documentclass{beamer}
\usepackage{ctex, hyperref}
\usepackage[T1]{fontenc}
%\usepackage{cite}
% other packages
\usepackage{latexsym,amsmath,xcolor,multicol,booktabs,calligra}
\usepackage{biblatex}
\usepackage{graphicx,pstricks,listings,stackengine}
\usefonttheme[onlymath]{serif}%统一字体
\addbibresource{ref.bib}
\renewcommand{\today}{\number\year 年\number\month 月\number\day 日}

%可以多个bib文件
\author[Dr. Zealer]{ Dr. Zealer }
\title[Why do you need to study some economics]{Why do you need to study some economics?}
\subtitle{——Let me tell you something}
\institute{河北经贸大学xx学院}
\usepackage{Tsinghua}

% defs下边框格式
\def\cmd#1{\texttt{\color{red}\footnotesize $\backslash$#1}}
\def\env#1{\texttt{\color{blue}\footnotesize #1}}
\definecolor{deepblue}{rgb}{0,0,0.5}
\definecolor{deepred}{rgb}{0.6,0,0}
\definecolor{deepgreen}{rgb}{0,0.5,0}
\definecolor{halfgray}{gray}{0.55}


\lstset{
    basicstyle=\ttfamily\small,
    keywordstyle=\bfseries\color{deepblue},
    emphstyle=\ttfamily\color{deepred},    % Custom highlighting style
    stringstyle=\color{deepgreen},
    numbers=left,
    numberstyle=\small\color{halfgray},
    rulesepcolor=\color{red!20!green!20!blue!20},
    frame=shadowbox,
}


\setbeamertemplate{bibliography item}[text]

\begin{document}

\kaishu
\begin{frame}
    \titlepage
    \begin{figure}[htpb]
        \begin{center}
            \includegraphics[width=0.25\linewidth]{logo2.jpg}
        \end{center}
    \end{figure}
\end{frame}

\begin{frame}
    \tableofcontents[sectionstyle=show,subsectionstyle=show/shaded/hide,subsubsectionstyle=show/shaded/hide]
\end{frame}





\section{研究背景}
\begin{frame}{背景与意义}
  
        
\end{frame}
\begin{frame}{国内外研究现状}
    
\end{frame}





\section{研究内容}
\begin{frame}{}
    
\end{frame}

\section{研究方法}
\begin{frame}{}
    
\end{frame}
\section{总结与展望}
\begin{frame}{总结}
    \begin{itemize}
        \item 
        \item 
        \item 
        \item 
    \end{itemize}
\end{frame}
\begin{frame}{不足之处}
   \cite{origin}
   %\cite{second}
\end{frame}

%\section{参考文献}

\begin{frame}

   %\cite{second}
\nocite{*}
\printbibliography[title={参考文献}]
%上述命令会生成一个新的frame,但不能加到目录，这一点暂未解决。
\end{frame}

\begin{frame}
    \begin{center}
        {\Huge\calligra Thanks!}
    \end{center}
\end{frame}

\end{document}

\end{document}